\documentclass[10pt]{extarticle}
\usepackage[a4paper,margin=1.5cm]{geometry}
\usepackage[T1]{fontenc}
\usepackage{amsmath,amssymb}
\usepackage{graphicx}
\usepackage[export]{adjustbox}
\usepackage{multicol}
\usepackage{enumitem}
\usepackage{booktabs,array,tabularx}
\setlength{\columnsep}{0.7cm}
\setlength{\parindent}{0pt}
\setlength{\parskip}{3pt}
\pagestyle{empty}
\makeatletter
\newcommand{\qfigdef}[4]{\@namedef{qf@file@#1}{#2}\@namedef{qf@w@#1}{#3}\@namedef{qf@wide@#1}{#4}}
\newcommand{\qf@check}[1]{\@ifundefined{qf@file@#1}{\PackageError{score2}{Unknown figure id #1}{}}{}}
\newcommand{\qf@img}[4][]{\includegraphics[width=#2\dimexpr\@nameuse{qf@w@#4}pt\relax,max width=#3,max height=0.45\textheight,#1]{\@nameuse{qf@file@#4}}}
\newif\ifqf@full
\newif\ifqf@boxed
\newsavebox{\qf@box}
\newcommand{\qf@open}{\global\qf@fulltrue\par\global\qf@boxedtrue
  \setbox\qf@box\vbox\bgroup\hsize\textwidth\linewidth\textwidth}
\newcommand{\qf@emit}[1]{\par\egroup\global\qf@boxedfalse
  \ifdim\dimexpr\ht\qf@box+\dp\qf@box\relax<#1\textheight\noindent\usebox{\qf@box}\par\else\unvbox\qf@box\fi}
\newcommand{\qfullsplit}{\ifqf@boxed\qf@emit{0.85}\fi}
\newcommand{\qf@close}{\ifqf@boxed\qf@emit{0.45}\fi\global\qf@fullfalse}
\newcommand{\qfullbegin}{\end{multicols}\qf@open}
\newcommand{\qfullend}{\qf@close\begin{multicols}{2}}
\newcommand{\qfullstart}{\qf@open}
\newcommand{\qfulldone}{\qf@close}
\newcommand{\qfig}[2][1]{\qf@check{#2}%
  \ifqf@full\qf@blk{#1}{#2}\else
  \ifnum\@nameuse{qf@wide@#2}=1 \expandafter\@firstoftwo\else\expandafter\@secondoftwo\fi
  {\end{multicols}\begin{center}\qf@img{#1}{\textwidth}{#2}\end{center}\begin{multicols}{2}}%
  {\qf@blk{#1}{#2}}\fi}
\newcommand{\qf@blk}[2]{\par\smallskip{\centering\qf@img{#1}{\linewidth}{#2}\par}\smallskip}
\newcommand{\qfigi}[2][1]{\qf@check{#2}\qf@img[valign=c]{#1}{\linewidth}{#2}}
\makeatother
\newcommand{\papersec}[1]{\par\bigskip{\centering\Large\bfseries PAPER-#1\par}\medskip}
\newcommand{\partsec}[2]{\par\medskip\noindent{\large\bfseries PART-#1: #2}\par\smallskip}
\newcommand{\sect}[1]{\par\smallskip\noindent\textbf{#1}\par}
\newcommand{\qn}[1]{\par\smallskip\noindent\textbf{#1.}~}
\newcommand{\paraq}[2]{\par\smallskip\noindent\textbf{Paragraph for Questions #1 and #2}\par}
\newenvironment{qoptions}{\begin{itemize}[leftmargin=1.8em,itemsep=1pt,topsep=2pt,label={}]}{\end{itemize}}
\newcommand{\opt}[2]{\item[(#1)] #2}
% figure registry (generated by score2_tools.py finalize)
\qfigdef{p1_phy_s1_q02_1}{SRG (Enthusiast) - Test 28 (20 Aug 2023)-img/p1_phy_s1_q02_1.png}{150.1}{0}
\qfigdef{p1_phy_s1_q03_1}{SRG (Enthusiast) - Test 28 (20 Aug 2023)-img/p1_phy_s1_q03_1.png}{245.2}{0}
\qfigdef{p1_phy_s1_q04_1}{SRG (Enthusiast) - Test 28 (20 Aug 2023)-img/p1_phy_s1_q04_1.png}{130.1}{0}
\qfigdef{p1_phy_s1_q05_1}{SRG (Enthusiast) - Test 28 (20 Aug 2023)-img/p1_phy_s1_q05_1.png}{71.5}{0}
\qfigdef{p1_phy_s1_q06_1}{SRG (Enthusiast) - Test 28 (20 Aug 2023)-img/p1_phy_s1_q06_1.png}{74.2}{0}
\qfigdef{p1_phy_s1_q08_1}{SRG (Enthusiast) - Test 28 (20 Aug 2023)-img/p1_phy_s1_q08_1.png}{94.8}{0}
\qfigdef{p1_phy_s2_q09_1}{SRG (Enthusiast) - Test 28 (20 Aug 2023)-img/p1_phy_s2_q09_1.png}{133.0}{0}
\qfigdef{p1_phy_s2_q10_1}{SRG (Enthusiast) - Test 28 (20 Aug 2023)-img/p1_phy_s2_q10_1.png}{137.3}{0}
\qfigdef{p1_phy_s3_q01_1}{SRG (Enthusiast) - Test 28 (20 Aug 2023)-img/p1_phy_s3_q01_1.png}{130.3}{0}
\qfigdef{p1_phy_s3_q03_1}{SRG (Enthusiast) - Test 28 (20 Aug 2023)-img/p1_phy_s3_q03_1.png}{159.4}{0}
\qfigdef{p1_phy_s3_q05_1}{SRG (Enthusiast) - Test 28 (20 Aug 2023)-img/p1_phy_s3_q05_1.png}{133.9}{0}
\qfigdef{p1_phy_s3_q06_1}{SRG (Enthusiast) - Test 28 (20 Aug 2023)-img/p1_phy_s3_q06_1.png}{175.8}{0}
\qfigdef{p1_phy_s3_q07_1}{SRG (Enthusiast) - Test 28 (20 Aug 2023)-img/p1_phy_s3_q07_1.png}{90.9}{0}
\qfigdef{p1_phy_s3_q08_1}{SRG (Enthusiast) - Test 28 (20 Aug 2023)-img/p1_phy_s3_q08_1.png}{88.6}{0}
\qfigdef{p1_phy_s3_q08_2}{SRG (Enthusiast) - Test 28 (20 Aug 2023)-img/p1_phy_s3_q08_2.png}{104.4}{0}
\qfigdef{p1_chem_s1_q05_1}{SRG (Enthusiast) - Test 28 (20 Aug 2023)-img/p1_chem_s1_q05_1.png}{152.3}{0}
\qfigdef{p1_chem_s3_q03_1}{SRG (Enthusiast) - Test 28 (20 Aug 2023)-img/p1_chem_s3_q03_1.png}{61.7}{0}
\qfigdef{p1_chem_s3_q05_1}{SRG (Enthusiast) - Test 28 (20 Aug 2023)-img/p1_chem_s3_q05_1.png}{226.8}{0}

\begin{document}
\begin{center}{\Large\bfseries SRG (Enthusiast) --- Test 28 (20 Aug 2023)}\end{center}
\begin{multicols}{2}
\partsec{1}{PHYSICS}
\sect{Section-I(i)}
\qn{1} A charged particle goes undeflected in a GRAVITY FREE region containing uniform electric and magnetic field. It is not possible at all ($\vec\upsilon$ = velocity of particle, $\vec E$ = uniform electric field, $\vec B$ = uniform magnetic field)
\begin{qoptions}
\opt{A}{$\vec E\,||\,\vec B,\ \vec\upsilon\,||\,\vec E$}
\opt{B}{$\vec E$ is not collinear to $\vec B$}
\opt{C}{$\vec\upsilon\,||\,\vec B$ but $\vec E$ is not collinear to $\vec B$}
\opt{D}{$\vec E\,||\,\vec B$ but $\vec\upsilon$ is not collinear to $\vec E$}
\end{qoptions}
\qn{2} A single circular loop of wire with radius 0.02 cm carries a current of 8.0 A. It is placed at the centre of a solenoid that has length 0.65 m, radius 0.080 m and 1300 turns. [you can treat solenoid as infinite]
\qfig{p1_phy_s1_q02_1}
\begin{qoptions}
\opt{A}{The value of the current in the solenoid so that the magnetic field at the centre of the loop becomes zero, is equal to 4.4 A.}
\opt{B}{The value of the current in the solenoid so that the magnetic field at the centre of the loop becomes zero, is equal to 10 A.}
\opt{C}{The magnitude of the total magnetic field at the centre of the loop (due to both the loop and the solenoid) if the current in the loop is reversed in direction from that needed to make the total field equal to zero Tesla, is $8\pi\times10^{-3}$T.}
\opt{D}{The magnitude of the total magnetic field at the centre of the loop (due to both the loop and the solenoid) if the current in the loop is reversed in direction from that needed to make the total field equal to zero Tesla, is $16\pi\times10^{-3}$T.}
\end{qoptions}
\qn{3} Figure shows cross section of two large parallel metal sheets carrying electric currents along their surfaces. The current in each sheet is $\dfrac{10}{\pi}$ A/m along the width. Consider two points A and B, as shown in the figure with their positions.
\qfig{p1_phy_s1_q03_1}
\begin{qoptions}
\opt{A}{Magnitude field at A is $4\mu$T along x-direction.}
\opt{B}{Magnetic field at A is $4\mu$T along negative x-direction.}
\opt{C}{Magnetic field at B is zero.}
\opt{D}{Magnetic field at B is $2\mu$T along x-direction.}
\end{qoptions}
\qn{4} The disc of radius $r$ is confined to roll without slipping at A and B. If the plates have the velocities as shown, then
\qfig{p1_phy_s1_q04_1}
\begin{qoptions}
\opt{A}{Linear velocity $v_0=v$}
\opt{B}{Angular velocity of disc is $\dfrac{3v}{2r}$}
\opt{C}{Angular velocity of disc is $\dfrac{2v}{r}$}
\opt{D}{None of these}
\end{qoptions}
\qn{5} A thin rod of mass $m$ and length $\ell$ is hinged to a ceiling and it is free to rotate in a vertical plane. A particle of mass $m$, moving with speed $v$ strikes it as shown in the figure and gets stick with the rod. Then
\qfig{p1_phy_s1_q05_1}
\begin{qoptions}
\opt{A}{The value of $v$, for which rod becomes horizontal after collision is $\sqrt{\dfrac{168}{9}g\ell}$}
\opt{B}{The value of $v$, for which rod becomes horizontal after collision is $\sqrt{\dfrac{53}{3}g\ell}$}
\opt{C}{Angular momentum of (rod + particle) system will remain constant about hinge just before and after collision}
\opt{D}{Angular momentum of (rod + particle) system will remain same about centre of mass just before and after collision}
\end{qoptions}
\qn{6} A thin uniform rod of mass $m$ and length $\ell$ is free to rotate about its upper end. When it is at rest, it receives an impulse $J$ at its lowest point, normal to its length. Immediately after impact:
\qfig{p1_phy_s1_q06_1}
\begin{qoptions}
\opt{A}{The angular momentum of the rod is $J\ell$.}
\opt{B}{The angular velocity of the rod is $\dfrac{3J}{m\ell}$}
\opt{C}{The kinetic energy of the rod is $\dfrac{3J^2}{2m}$}
\opt{D}{The linear velocity of the midpoint of the rod is $\dfrac{3J}{2m}$}
\end{qoptions}
\qn{7} A bicycle is in motion. The force of friction exerted by the ground on its wheel is such that it acts:
\begin{qoptions}
\opt{A}{In backward direction on front wheel and in forward direction on rear wheel when it is accelerating}
\opt{B}{In forward direction on front wheel and in backward direction on rear wheel when brakes are applied on rear wheel only}
\opt{C}{In backward direction on front wheel and in forward direction on rear wheel when brakes are applied on rear wheel only}
\opt{D}{In backward direction on both the wheels when brakes are applied on front wheel}
\end{qoptions}
\qn{8} If a charged particle of charge to mass ratio $\dfrac qm=\alpha$ is entering in a magnetic field of strength $B$ perpendicular only at a speed $v=(2\alpha d)(B)$, then which of the following is correct?
\qfig{p1_phy_s1_q08_1}
\begin{qoptions}
\opt{A}{Angle subtended by charged particle at the centre of circular path is $2\pi$}
\opt{B}{The charge will move on a circular path and will come out from magnetic field at a distance $4d$ from the point of insertion}
\opt{C}{The time for which particle will be in the magnetic field is $\dfrac{2\pi}{\alpha B}$}
\opt{D}{The charged particle will subtend an angle of $90^\circ$ at the centre of circular path}
\end{qoptions}
\sect{Section-I(ii)}
\qn{9} A disc of radius $R$ is rolling without slipping with an angular acceleration $\alpha$, on a fixed horizontal plane. Four points are marked at the end of horizontal and vertical diameter of a circle $r\,(<R)$ on the disc. If horizontal vertical direction are chosen as $x$ and $y$ axis as shown in the figure, then acceleration of points 1, 2, 3 and 4 are $\vec a_1,\vec a_2,\vec a_3$ and $\vec a_4$ respectively, at the moment when angular velocity of the disc is $\omega$. Match the following
\qfig{p1_phy_s2_q09_1}
\begin{center}\small\renewcommand{\arraystretch}{1.4}
\begin{tabular}{|c|l||c|l|}\hline
\multicolumn{2}{|c||}{Column - I} & \multicolumn{2}{c|}{Column - II}\\\hline
A & $\vec a_1$ & P & $(R\alpha-r\alpha)\hat i+(r\omega^2)\hat j$\\\hline
B & $\vec a_2$ & Q & $(R\alpha+r\alpha)\hat i-(r\omega^2)\hat j$\\\hline
C & $\vec a_3$ & R & $(R\alpha-r\omega^2)\hat i-(r\alpha)\hat j$\\\hline
D & $\vec a_4$ & S & $(R\alpha+r\omega^2)\hat i+(r\alpha)\hat j$\\\hline
 & & T & None of these\\\hline
\end{tabular}
\end{center}
\begin{qoptions}
\opt{A}{$A\to Q,\ B\to R,\ C\to P,\ D\to S$}
\opt{B}{$A\to Q,\ B\to P,\ C\to R,\ D\to S$}
\opt{C}{$A\to Q,\ B\to P,\ C\to S,\ D\to R$}
\opt{D}{$A\to R,\ B\to P,\ C\to S,\ D\to P$}
\end{qoptions}
\qn{10} A very small current carrying square loop (current $I$) of side `$L$' is placed in $y$-$z$ plane with centre at origin of the coordinate system (shown in figure). In column-I the coordinate of the points are given and in column-II magnitude of strength of magnetic field is given. Then
\qfig{p1_phy_s2_q10_1}
\begin{center}\small\renewcommand{\arraystretch}{2.6}
\begin{tabular}{|c|l||c|c|}\hline
\multicolumn{2}{|c||}{Column - I} & \multicolumn{2}{c|}{Column - II}\\\hline
A & At point $O\,(0,0,0)$ & P & $\dfrac{\mu_0IL^2}{2\pi a^3}$\\\hline
B & At point $P_1\,(a,0,0)$ (here $a\gg L$) & Q & $\dfrac{2\sqrt2\,\mu_0I}{\pi L}$\\\hline
C & At point $P_2\,(0,a,0)$ (here $a\gg L$) & R & $\dfrac{\mu_0\sqrt5\,IL^2}{16\pi a^3}$\\\hline
D & At point $P_3\,(a,a,0)$ (here $a\gg L$) & S & $\dfrac{\mu_0IL^2}{4\pi a^3}$\\\hline
 & & T & $\dfrac{\mu_0\sqrt5\,IL^2}{4\pi a^3}$\\\hline
\end{tabular}
\end{center}
\begin{qoptions}
\opt{A}{$A\to P,\ B\to Q,\ C\to S,\ D\to R$}
\opt{B}{$A\to Q,\ B\to P,\ C\to S,\ D\to R$}
\opt{C}{$A\to P,\ B\to Q,\ C\to R,\ D\to S$}
\opt{D}{$A\to P,\ B\to S,\ C\to R,\ D\to Q$}
\end{qoptions}
\sect{Section-II}
\qn{1} A uniform rod of length $R$ and mass $M$ is free to rotate about a horizontal axis passing through hinge $P$ as in figure. First it is taken aside such that it becomes horizontal and then released. At the lowest point the rod hits the block $B$ of mass $m$ and stops. If the ratio of mass of rod $M$ to the mass of block $m$ such that the block $B$ completes the circle is $\sqrt x$ then $x$ is. Neglect any friction.
\qfig{p1_phy_s3_q01_1}
\qn{2} An electron (mass $=9\times10^{-31}$ kg, charge $=-1.6\times10^{-19}$C) experiences no deflection if subjected to an electric field of $3.2\times10^5$ V/m and a magnetic field of $2.0\times10^{-3}$ Wb/m$^2$. Both the fields are normal to the path of electron and to each other. If the electric field is removed, then the electron will revolve in an orbit of radius \_\_\_\_\_\_\_ meters. [assume gravity free space]
\qn{3} A conducting ring of mass 2 kg and radius 0.5 m is placed on a smooth horizontal plane. The ring carries a current $i=4$A. A horizontal magnetic field $B=10$ T is switched on at time as shown in figure. The initial angular acceleration of the ring will be \_\_\_\_ $\pi$rad/s$^2$
\qfig{p1_phy_s3_q03_1}
\qn{4} A moving coil galvanometer has 150 equal divisions. Its current sensitivity is 10 divisions per milli ampere and voltage sensitivity is 2 divisions per millivolt. In order that each division reads 1V, the resistance in Ohm's needed to be connected in series with the coil will be \_\_\_\_\_\_ $\Omega$
\qn{5} A conducting wire bent in the form of a parabola $y^2=2x$ carries a current $i=2$A as shown in figure. This wire is placed in a uniform magnetic field $\vec B=-4\hat k$ Telsa. The magnitude of magnetic force on the wire is \_\_\_\_ (in newton)
\qfig{p1_phy_s3_q05_1}
\qn{6} A uniform rod of length 75 cm is hinged at one of its ends and is free to rotate in vertical plane. It is released from rest when rod is horizontal. When the rod becomes vertical, it is broken at mid-point without any jerk and lower part now moves freely. The distance of centre of lower part from hinge, when it again becomes vertical for the first time is $r$. Find the approximate value of $2r$ in meters
\qfig{p1_phy_s3_q06_1}
\qn{7} A thin uniform square plate ABCD of side `$a$' and mass $m$ is suspended in vertical plane as shown in the figure. AE and BF are two massless inextensible strings. The line AB is horizontal. Find the tension in the string AE just after BF is cut is $\dfrac{2mg}{x}$ then $x$ is.
\qfig{p1_phy_s3_q07_1}
\qn{8} A carpet of mass $M$ made of inextensible material is rolled along its length in the form of a cylinder of radius $R$ and is kept on a rough floor. The carpet starts unrolling without sliding on the floor when a negligibly small push is given to it. The horizontal velocity of the axis of the cylindrical part of the carpet when its radius reduces to $\dfrac R2$. Is $\sqrt{\dfrac{xgR}{3}}\ \text{ms}^{-1}$ then $x$ is:
\qfig{p1_phy_s3_q08_1}
\qfig{p1_phy_s3_q08_2}
\partsec{2}{CHEMISTRY}
\sect{Section-I(i)}
\qn{1} Identify the correct statement(s) regarding hydrogen gas
\begin{qoptions}
\opt{A}{Unlike $\mathrm{H_2}$, $\mathrm{D_2}$ do not show Nuclear spin isomerism.}
\opt{B}{Orthohydrogen differs from parahydrogen w.r.t their b.p., sp. Heat and thermal conductivity}
\opt{C}{As the temperature is increased from absolute zero some of the para form gradually changes to ortho form.}
\opt{D}{Parahydrogen has lower internal energy than orthohydrogen.}
\end{qoptions}
\qn{2} Which one of the following processes will produce hard water?
\begin{qoptions}
\opt{A}{Saturation of water with $\mathrm{Ca(HCO_3)_2}$ solution}
\opt{B}{Saturation of water with $\mathrm{MgCO_3}$}
\opt{C}{Saturation of water with $\mathrm{CaSO_4}$}
\opt{D}{Addition of $\mathrm{Na_3PO_4}$ to water}
\end{qoptions}
\qn{3} Which of the following is correct?
\begin{qoptions}
\opt{A}{Oxidation state of Fe in sodium Nitro prusside is $+2$}
\opt{B}{AgI is in yellow color due to polarization effect}
\opt{C}{AgBr is a pale yellow ppt. and is sparingly soluble in excess dil. $\mathrm{NH_4OH}$}
\opt{D}{Yellow ppt is the detection of phosphorus when an organic compound is heated with $\mathrm{Na_2O_2}$ and then boiled with con. $\mathrm{HNO_3}$, followed by the addition of Ammonium molybdate is $\mathrm{(NH_4)_3PO_4.12MoO_3}$.}
\end{qoptions}
\qn{4} N, S, halogens and phosphorous present in an organic compound are detected by Lassaigne's test, by fusing the compound with Na metal. The compounds formed are extracted from fused by boiling it with distilled water. This extract is known as sodium fusion extract or Lassaigne's extract.\\
An organic compound containing N, S, and O as extra elements is fused with metallic sodium and then extracted with distilled water. Which species is / are not present in the Lassaigne's extract?
\begin{qoptions}
\opt{A}{$\mathrm{NO_3^-}$}
\opt{B}{$\mathrm{CN^-}$}
\opt{C}{$\mathrm{SCN^-}$}
\opt{D}{$\mathrm{S^{2-}}$}
\end{qoptions}
\qn{5} The correct statement(s) about the following sugars $\mathbf X$ and $\mathbf Y$ is (are)
\qfig{p1_chem_s1_q05_1}
\begin{qoptions}
\opt{A}{$\mathbf X$ is a reducing sugar and $\mathbf Y$ is a non-reducing sugar.}
\opt{B}{$\mathbf X$ is a non-reducing sugar and $\mathbf Y$ is a reducing sugar.}
\opt{C}{The glucosidic linkages in $\mathbf X$ and $\mathbf Y$ are $\alpha$ and $\beta$, respectively}
\opt{D}{The glucosidic linkages in $\mathbf X$ and $\mathbf Y$ are $\beta$ and $\alpha$, respectively}
\end{qoptions}
\qn{6} Nylon is used as a general name for all synthetic fibres forming polyamides. Identify the correct statement about Nylon.
\begin{qoptions}
\opt{A}{Nylon-6,10 is a condensation polymerisation of hexamethylenediamine and sebacic acid.}
\opt{B}{Nylon-2-Nylon-6 is alternating polyamide copolymer of glycine and amino caproic acid.}
\opt{C}{Nylon-6 is produced by prolonged heating of caprolactum with $\mathrm{H_2O}/\Delta$.}
\opt{D}{Nylon-6 on further polymerisation gives Nylon-6,6.}
\end{qoptions}
\qn{7} Neoprene, the first synthetic rubber manufactured on large scale, has some advantages over natural rubber. Identify \textbf{CORRECT} statement(s) from the following options.
\begin{qoptions}
\opt{A}{It is resistant to sunlight, oxygen, ozone and heat.}
\opt{B}{It is non-inflammable.}
\opt{C}{Monomer of this polymer, can be prepared by using $\mathrm{C_2H_2}$.}
\opt{D}{It is resistant to actions of oils.}
\end{qoptions}
\qn{8} A cubical crystal is such that the corners atoms are occupied by P atoms, face centers are occupied by Q atoms, alternate tetrahedral voids are occupied by R atoms and edge centers are occupied by S atoms. Based on this info identify the correct option.
\begin{qoptions}
\opt{A}{If there are no defects in the crystal then the formula should be $\mathrm{PQ_3R_4S_3}$.}
\opt{B}{If all the atoms lying on one of the body diagonal are missing then formula should be $\mathrm{PQ_4R_{8/3}S_4}$}
\opt{C}{If all the atoms lying on a line passing connecting edge centers are missing then formula should be $\mathrm{P_2Q_5R_8S_5}$}
\opt{D}{If all the atoms lying on a line passing through body center \& connecting opposite face centers are missing then formula will be $\mathrm{PQ_2R_4S_3}$.}
\end{qoptions}
\sect{Section-I(ii)}
\qn{9} The correct match between item I and item II is
\begin{center}\small
\begin{tabular}{|l|p{0.4\linewidth}|}\hline
\textbf{Item I} & \textbf{Item II}\\\hline
(a) High-density polythene & (I) Peroxide catalyst\\\hline
(b) Polyacrylonitrile & (II) Condensation at high temperature and pressure\\\hline
(c) Novolac & (III) Zeigler-Natta catalyst\\\hline
(d) Nylon-6 & (IV) Acid or base catalyst\\\hline
\end{tabular}
\end{center}
\begin{qoptions}
\opt{A}{(a)-III, (b)-I, (c)-IV, (d)-II}
\opt{B}{(a)-I, (b)-III, (c)-II, (d)-IV}
\opt{C}{(a)-II, (b)-I, (c)-III, (d)-IV}
\opt{D}{(a)-IV, (b)-III, (c)-I, (d)-II}
\end{qoptions}
\qn{10} Match the medicines given in column I with their use given in column II and select the correct option:
\begin{center}\small
\begin{tabular}{|l|l||l|l|}\hline
\multicolumn{2}{|c||}{\textbf{Column - I}} & \multicolumn{2}{c|}{\textbf{Column - II}}\\\hline
(A) & Ranitidine & (i) & Tranquilizer\\\hline
(B) & Furacine & (ii) & Antibiotic\\\hline
(C) & Phenelzine & (iii) & Antihistamine\\\hline
(D) & Salvarsan & (iv) & Antiseptic\\\hline
\end{tabular}
\end{center}
\begin{qoptions}
\opt{A}{(A)-(iii), (B)-(iv), (C)-(i), (D)-(ii)}
\opt{B}{(A)-(iv), (B)-(i), (C)-(ii), (D)-(iii)}
\opt{C}{(A)-(i), (B)-(ii), (C)-(iii), (D)-(iv)}
\opt{D}{(A)-(ii), (B)-(iii), (C)-(iv), (D)-(i)}
\end{qoptions}
\sect{Section-II}
\qn{1} In how many of the following reactions hydrogen gas is produced as one of the product.\\
i) Electrolysis of acidified water.\\
ii) Passing steam over red hot coke.\\
iii) Reaction of aluminium with excess of sodium hydroxide.\\
iv) Reaction of Zinc with dil. HCl.\\
v) Pyrolysis of hydrocarbons.\\
vi) Reaction of sodium amalgam with water.\\
vii) By dissolving hydrolith in water.\\
viii) By heating a mixture of steam and hydro carbons.\\
ix) Electrolysis of very very diluted NaCl aq.solution.
\qn{2} During the estimation of nitrogen in an organic compound by Kjeldhal's method, the ammonia evolved from 0.5 g of the compound, was absorbed in 50 ml of 0.5 M $\mathrm{H_2SO_4}$. The residual acid required 60 ml of 0.5 M NaOH solution. The percentage of nitrogen in the compound is \_\_\_\_\_\_
\qn{3} What is the packing fraction of the two dimensional unit cell shown in figure?
\qfig{p1_chem_s3_q03_1}
\qn{4} Density of one form of $\mathrm{CaCO_3}$ solid is 2.93 g/cc and it has orthorhombic unit cell with unit cell parameter: $a=4.6$\,\AA, $b=5.7$\,\AA\ and $c=57$\,\AA. Calculate number of $\mathrm{Ca^{2+}}$ ions per unit cell.\\
{[}Give the answer to closest integer]
\qn{5} Identify the total number of carbohydrates which give mutarotation
\qfig{p1_chem_s3_q05_1}
\qn{6} 0.75 g of an ester of a monocarboxylic acid is refluxed with 100 ml of 0.125 M KOH and the excess KOH neutralised by back titration with 27.5 ml of 0.3 M HCl. What is its saponification value?
\qn{7} The $\mathrm{pK_a}$ of the carboxyl group in an amino acid valine, $\mathrm{(CH_3)_2CHCH(NH_2)(COOH)}$ is 2.31 and the $\mathrm{pK_a}$ for the amino group of the same amino acid is 9.69. Compute the isoelectric point (pI) for valine.
\qn{8} In E.coli DNA, the AT/GC ratio is 0.93. If the number of moles of adenine in its DNA sample is 465000, The number of guanine unit is $x\times10^{y}$. What is the value of $x+y$. \{Both $x$ \& $y$ are integer\}
\partsec{3}{MATHEMATICS}
\sect{Section-I(i)}
\qn{1} Which of the following are true:
\begin{qoptions}
\opt{A}{$\sqrt{e^\pi}>\left(\dfrac\pi2\right)^e$}
\opt{B}{$e^\pi>\pi^e$}
\opt{C}{$e^3>3^e$}
\opt{D}{$e^2>2^e$}
\end{qoptions}
\qn{2} Interval to which $k$ may belong so that function $g(x)=\left(1-\dfrac{\sqrt{21-4k-k^2}}{k+1}\right)x^3+5x+\sqrt6$ is increasing at every point of its domain is:
\begin{qoptions}
\opt{A}{$(2,3]$}
\opt{B}{$[-7,-1)$}
\opt{C}{$[-6,-2]$}
\opt{D}{$(2,2.5]$}
\end{qoptions}
\qn{3} There exists a triangle ABC satisfying the conditions
\begin{qoptions}
\opt{A}{$b\sin A=a,\ A<\dfrac\pi2$}
\opt{B}{$b\sin A>a,\ A>\dfrac\pi2$}
\opt{C}{$b\sin A>a,\ A<\dfrac\pi2$}
\opt{D}{$b\sin A<a,\ A<\dfrac\pi2,\ b>a$}
\end{qoptions}
\qn{4} $f:[0,2]\to R$ is continuous on $[0,2]$ and twice differentiable in $(0,2)$. If $f(0)=0$, $f(1)=1$, $f(2)=1$ then which options are correct.
\begin{qoptions}
\opt{A}{$2f'(c)+2c=3$ for some $c\in(0,1)$}
\opt{B}{$2f'(c)+2c=3$ for some $c\in(1,2)$}
\opt{C}{$f''(c)=-1$ for some $c\in(0,2)$}
\opt{D}{$f'(c)=0$ for some $c\in(1,2)$}
\end{qoptions}
\qn{5} In a triangle with sides $a$, $b$ and $c$, a semicircle touching the sides AC and CB is inscribed whose diameter lies on AB. Then, the radius of the semicircle is
\begin{qoptions}
\opt{A}{$a/2$}
\opt{B}{$\Delta/s$}
\opt{C}{$\dfrac{2\Delta}{a+b}$}
\opt{D}{$\dfrac{2abc}{(s)(a+b)}\cos\dfrac A2\cos\dfrac B2\cos\dfrac C2$}
\end{qoptions}
\qn{6} If $x^3-2x^2y^2+5x+y-5=0$ and $y(1)=1$, then
\begin{qoptions}
\opt{A}{$y'(1)=4/3$}
\opt{B}{$y''(1)=-4/3$}
\opt{C}{$y''(1)=-\dfrac{238}{27}$}
\opt{D}{$y'(1)=2/3$}
\end{qoptions}
\qn{7} If $y=\dfrac{x^4-x^2+1}{x^2+\sqrt3x+1}$ and $\dfrac{dy}{dx}=ax+b$, then
\begin{qoptions}
\opt{A}{$a=\sqrt3$}
\opt{B}{$b=-\sqrt3$}
\opt{C}{$a-b=\cot\dfrac\pi{12}$}
\opt{D}{$a+b=\tan\dfrac\pi3$}
\end{qoptions}
\qn{8} If $y=\cos^{-1}\left(\dfrac{2x}{1+x^2}\right)$, then $\dfrac{dy}{dx}$ is
\begin{qoptions}
\opt{A}{$\dfrac{-2}{1+x^2}$ for all $x$}
\opt{B}{$\dfrac{-2}{1+x^2}$ for all $|x|<1$}
\opt{C}{$\dfrac{2}{1+x^2}$ for all $|x|>1$}
\opt{D}{None of these}
\end{qoptions}
\sect{Section-I(ii)}
\qn{9} Match the following
\begin{center}\small\renewcommand{\arraystretch}{1.3}
\begin{tabular}{|c|p{0.52\linewidth}|c|c|}\hline
\multicolumn{2}{|c|}{\textbf{Column-I}} & \multicolumn{2}{c|}{\textbf{Column-II}}\\\hline
(A) & In a triangle ABC, if $a^4-2\left(b^2+c^2\right)a^2+b^4+b^2c^2+c^4=0$, then $\angle A$ is can be & (P) & $30^\circ$\\\hline
(B) & In a triangle ABC, If $a^4+b^4+c^4=a^2b^2+2b^2c^2+2c^2a^2$, then $\angle C$ is can be & (Q) & $60^\circ$\\\hline
(C) & In a triangle ABC, If $a^4+b^4+c^4+2a^2c^2=2a^2b^2+2b^2c^2$, then $\angle B$ is & (R) & $90^\circ$\\\hline
(D) & In a triangle ABC, If $\dfrac1{a+c}+\dfrac1{b+c}=\dfrac3{a+b+c}$, then $\angle C$ is & (S) & $135^\circ$\\\hline
\end{tabular}
\end{center}
\begin{qoptions}
\opt{A}{(A) $\to$ (Q), (B) $\to$ (S), (C) $\to$ (P), (D) $\to$ (Q).}
\opt{B}{(A) $\to$ (R), (B) $\to$ (P), (C) $\to$ (S), (D) $\to$ (Q).}
\opt{C}{(A) $\to$ (Q), (B) $\to$ (P), (C) $\to$ (R), (D) $\to$ (Q).}
\opt{D}{(A) $\to$ (P), (B) $\to$ (R), (C) $\to$ (Q), (D) $\to$ (S).}
\end{qoptions}
\qn{10} Consider $f(x)=\left|\ln|x|\right|-kx^2,\ k\ne0,\ x\ne0$. Match the column I with the value of $k$ in column II.
\begin{center}\small\renewcommand{\arraystretch}{1.8}
\begin{tabular}{|c|p{0.3\linewidth}||c|l|}\hline
\multicolumn{2}{|c||}{\textbf{Column-I}} & \multicolumn{2}{c|}{\textbf{Column-II}}\\\hline
(A) & $f(x)=0$ has two distinct solutions & (P) & $k\in(-\infty,\,0)$\\\hline
(B) & $f(x)=0$ has four distinct solutions & (Q) & $k=\dfrac1{2e}$\\\hline
(C) & $f(x)=0$ has six distinct solutions & (R) & $k\in\left(\dfrac1{2e},\infty\right)$\\\hline
(D) & $f(x)=0$ has no solution & (S) & $k\in\left(0,\dfrac1{2e}\right)$\\\hline
\end{tabular}
\end{center}
\begin{qoptions}
\opt{A}{(A) $\to$ (Q), (B) $\to$ (S), (C) $\to$ (R), (D) $\to$ (P).}
\opt{B}{(A) $\to$ (Q), (B) $\to$ (S), (C) $\to$ (P), (D) $\to$ (R).}
\opt{C}{(A) $\to$ (S), (B) $\to$ (Q), (C) $\to$ (R), (D) $\to$ (P).}
\opt{D}{(A) $\to$ (R), (B) $\to$ (Q), (C) $\to$ (S), (D) $\to$ (P).}
\end{qoptions}
\sect{Section-II}
\qn{1} In triangle ABC, $a=5$, $b=4$, $c=3$, $G$ is the centroid of triangle. If $R_1$ be the circumradius of triangle GAB then the value of $\dfrac{3R_1}{\sqrt{13}}$ is
\qn{2} Find total number of points of maxima and minima of the function\\
$f(x)=\begin{cases}x^4-2x^2; & -3<x\le0\\ 3x^{\frac23}-2x; & 0<x<2\\ -1; & x=0\\ |x-3|; & 2<x<4\end{cases}$
\qn{3} Let $f(x)=x^3+x^2+2x-1$. The minimum value of $[x]$ that satisfy $f(f(x))>f(2x+1)$ is (where $[x]$ denotes the greatest integer function)
\qn{4} If $f\!\left(\dfrac{x+y}{3}\right)=\dfrac{2+f(x)+f(y)}{3}$ for all real $x$ and $y$ and $f'(2)=2$, then $f(3)$ is equal to
\qn{5} A cylindrical container is to be made from certain solid material with the following constraints: It has a fixed inner volume of $V$ mm$^3$, has a 2 mm thick solid wall and is open at the top. The bottom of the container is a solid circular disc of thickness 2 mm and is of radius equal to the outer radius of the container.\\
If the volume of the material used to make the container is minimum when the inner radius of the container is 10 mm, then the value of $\dfrac V{1000}$ is \_\_\_\_\_\_\_.
\qn{6} Let $\Delta(p)=\begin{vmatrix}\cos p & p^3 & p^5\\ \sec^3p & \sin p & 0\\ \sin^2p & -p & p+p^3\end{vmatrix}$. The value of $\left.\dfrac{d^3(\Delta(p))}{dp^3}\right|_{p=0}$ is equal to
\qn{7} For a twice differentiable function $f(x)$, $g(x)$ is defined as $g(x)=\left(f'(x)\right)^2+f(x)\cdot f''(x)$, on $(a,e)$. If $a<b<c<d<e$, $f(a)=0$, $f(b)=2$, $f(c)=-1$, $f(d)=2$, $f(e)=0$ then minimum number of roots of the equation $g(x)=0$ in $(a,e)$ is
\qn{8} Let $f(\theta)=\sin\left(\tan^{-1}\left(\dfrac{\sin\theta}{\sqrt{\cos2\theta}}\right)\right)$, where $-\dfrac\pi4<\theta<\dfrac\pi4$. Then the value of $\dfrac{d}{d(\tan\theta)}(f(\theta))$ is
\end{multicols}
\end{document}
